\documentclass[10pt]{article}
\usepackage[utf8]{inputenc}
\usepackage[T1]{fontenc}
\usepackage{amsmath}
\usepackage{amsfonts}
\usepackage{amssymb}
\usepackage[version=4]{mhchem}
\usepackage{stmaryrd}
\usepackage{graphicx}
\usepackage[export]{adjustbox}
\graphicspath{ {./images/} }
\usepackage{caption}

\begin{document}
\captionsetup{singlelinecheck=false}
\section*{4- §. Grin funksiyasi usuli}
Elliptik tipdagi tenglamalar uchun chegaraviy masalalarni yechishda (Manba funksiya) Grin funksiyasi muhim oˋrin tutadi. Boˋlaklari silliq $S$ sirt bilan chegaralangan $D$ sohani qaraylik.\\
4.1-ta'rif. (4.2) Laplas tenglamasi uchun $D$ sohada ichki (tashqi) Dirixle masalasining Grin funksiyasi deb, ikkita $x, \xi \in D \cup S$ nuqtalarning funksiyasi boˋlgan va quyidagi shartlarni qanoatlantiruvchi $G(x, \xi)$ funksiyaga aytiladi [15]:

\begin{enumerate}
  \item $G(x, \zeta)$ funksiya ushbu ko rinishga ega
\end{enumerate}


\begin{equation*}
G(x, \xi)=\frac{1}{4 \pi}[E(x, \xi)+g(x, \xi)], \tag{4.56}
\end{equation*}


bu yerda $E(x, \xi)$ - Laplas tenglamasining (4.9) formula bilan aniqlangan fundamental (elementar) yechimi, $g(x, \xi)$ - funksiya esa $x \in D$ bo yicha ham, $\xi \in D$ bo yicha ham garmonik funksiya bo lib, $x$ bo yicha har bir $\xi$ lar uchun $\bar{D}$ sohada uzluksiz.\\
2) $x, \xi$ nuqtalarning hech bo Imaganda bittasi $S$ ga tegishli bo'lsa,


\begin{equation*}
G(x, \xi)=0 \tag{4.57}
\end{equation*}


bo'ladi. Agar $D$ soha chegaralanmagan bo'lsa, u holda $|x| \rightarrow \infty$ da $g(x, \xi) \rightarrow 0$ shart bajarilishi talab qilinadi.

Bu ta'rifga asosan, $G(x, \xi)$ funksiya $\xi$ nuqtadan tashqari barcha $D$ sohada garmonik bo lib, $g(x, \xi)$ funksiya yordamida aniqlanadi. $g(x, \xi)$ funksiya esa o'z navbatida $D$ sohada garmonik bo'lib, $S$ da $-\ln (1 / r), n=2$ yoki $-r^{2-n}, n>2$ qiymatlarga teng. Bu yerda $g(x, \xi)$ shunday garmonik funksiyaki, u chegarada maxsus qiymatlarni qabul qiladi. Ayrim hollarda bunday funksiyani topish ancha qulay bo'ladi. Bu funksiya maˋlum boˋlgandan so'ng, chegarada ixtiyoriy qiymatni qabul qiluvchi garmonik funksiyani topish mumkin bo ladi.

\section*{Grin funksiyasining xossalari:}
\begin{enumerate}
  \item $D \cup S$ da $G(x, \xi) \geq 0$;
  \item agar $D$ sohaning $S$ chegarasi yetarlicha silliq sirt bo lsa, u holda $G(x, \xi)$ funksiya $S$ chegarada har bir $\xi \in D$ da to'gˋri normal $\frac{\partial G(x, \xi)}{\partial n_{x}}$ hosilaga ega;
  \item $G(x, \xi$ funksiya $x \in D$ va $\xi \in D$ nuqtalarga nisbatan simmetrik funksiyadir, ya'ni $\quad G(x, \xi)=G(\xi, x)$;
  \item $x \in D$ nuqtalarda quyidagi tenglik
\end{enumerate}

$$
-\frac{1}{(n-2)\left|S_{1}\right|} \int_{S} \frac{\partial G(x, \xi)}{\partial n_{\xi}} d_{\xi} S=1
$$

oˋrinli. $\left|S_{1}\right|$ - birlik sferaning yuzi.\\
Agar $w=w(z)$ funksiya $D(n=2$ bo'lganda sohani birlik doiraga (ya'ni $|x|<1 g a$ ) konform akslantiruvchi funksiya boˋlsa, u holda Grin funksiyasi quyidagi formuladan

$$
G(z, \zeta)=\frac{1}{2 \pi} \ln \frac{1}{|w(z, \zeta)|}, \quad w(z, \zeta)=\frac{w(z)-w(\zeta)}{1-w(z) \cdot w(\zeta)},
$$

$z=x+i y, \zeta=\xi+i \eta$ topiladi.\\[0pt]
Grin funksiyasini tuzish. Qaralayotgan ko pgina sohalarda Grin funksiyasi akslantirsh usulidan foydalanib tuziladi [5], [10], [19].

I- masala. (4.2) Laplas tenglamasi uchun sharda ichki Dirixle masalasining Grin funksiyasini tuzing.\\
Yechish. Akslantirish usuli yordamida sharda Grin funksiyasini tuzamiz. $\mathrm{U}_{R}$ soha markazi koordinata boshida va radiusi $R$ ga teng bo"lgan shar bo"lsin. Uni chegaralab turgan sferani $S_{R}$ orqali belgilab olamiz. Faraz qilaylik, $y \neq 0$ bo lib, $y$ ga $S_{R}$ sferaga nisbatan simmetrik nuqta $y^{*}$ bo' $1 \sin$, ya'ni


\begin{equation*}
\left|y \|\left|y^{\cdot}\right|=R^{2} .\right. \tag{4.58}
\end{equation*}


$|y|<R$ bo'lganligi uchun, $y^{*}$ nuqta $S_{R}$ sferadan tashqarida yotadi. Bu nuqta uchun ushbu inversiya almashtirish o'rinlidir:

$$
y=\frac{R^{2}}{\left|y^{*}\right|^{2}} y^{*}, \quad y^{*}=\frac{R^{2}}{|y|^{2}} y, \quad|y|=\sqrt{y_{1}^{2}+y_{2}^{2}+y_{3}^{2}}
$$

yoki


\begin{equation*}
y_{i}^{*}=\frac{R^{2}}{|y|^{2}} y_{i}, \quad y_{i}=\frac{R^{2}}{\left|y^{*}\right|^{2}} y_{i}^{*}, \quad i=1,2,3 . \tag{4.59}
\end{equation*}


Grin funksiyani quyidagi ko'rinishda izlaymiz:


\begin{equation*}
G(x, y)=v(x, y)+g(x, y), \tag{4.60}
\end{equation*}


bu yerda $v(x, y)=1 / 4 \pi|x-y| \quad$ funksiya Laplas tenglamasining fundamental yechimi bo lib, bu funksiya $y \neq x$ bo lganda $y$ bo yicha ham, $x$ bo yicha ham Laplas tenglamasini qanoatlantiradi. $g(x, y)=-\frac{\alpha}{4 \pi\left|x-y^{*}\right|} \quad$ funksiya esa $\mathrm{U}_{k}$ da $y$ bo'yicha ham, $x$ bo yicha ham garmonik funksiya bo lib, $C^{\infty}\left(\mathrm{U}_{\mathrm{R}}\right)$ sinfga tegishlidir, ya'ni bu funksiyaning maxraji nolga aylanmaydi, $\alpha$ - keyinchalik aniqlanadigan son.

\begin{figure}[h]
\begin{center}
  \includegraphics[alt={},max width=\textwidth]{f9c87acf-24e6-44a5-be24-c530e7441594-03_326_511_512_365}
\captionsetup{labelformat=empty}
\caption{4.1 - chizma}
\end{center}
\end{figure}

Agar $x \in S_{R}$ bo'lsa, u holda 4.1 - chizmaga ko'ra $O x y{ }^{*}$ va $O x y$ uchburchaklar o'xshash bo'ladi, chunki ular umumiy $O$ burchakka ega va bu burchakni hosil qilgan tomonlari (4.58) va (4.59) tengliklarga koˋra proporsionaldir, ya'ni

$$
\frac{\left|y^{*}\right|}{R}=\frac{R}{y},
$$

bu yerda $\left|y^{*}\right|$ va $R$ - $O x y^{*}$ va $O x y$ uchburchaklarning mos ravishda tomonlarining uzunliklari.

Shunday qilib, uchburchaklarning o'xshashligidan


\begin{equation*}
R \cdot|x-y|=|y| \cdot\left|x-y^{*}\right| \quad \text { yoki } \quad R \cdot|x-y|=\left|x-y^{*}\right| \cdot|y| \tag{4.61}
\end{equation*}


tengliklarni hosil qilamiz. Grin funksiyasining ta’rifiga ko’ra $\alpha$ sonini shunday tanlaymizki, $\left.G(x, y)\right|_{S_{R}}=0$ boˋlsin. Bundan va (4.60) dan


\begin{equation*}
\alpha \cdot|x-y|=\left|x-y^{*}\right| \Rightarrow \alpha=\frac{\left|x-y^{*}\right|}{|x-y|}, \quad x \in S_{R} \tag{4.62}
\end{equation*}


kelib chiqadi. (4.61) ga asosan esa


\begin{equation*}
\alpha=R / y \tag{4.63}
\end{equation*}


ega bo"lamiz.\\
Shunday qilib, (4.60) va (4.63)ga koˋra


\begin{equation*}
G(x, y)=\frac{1}{4 \pi|x-y|}-\frac{R}{4 \pi|y|\left|x-y^{0}\right|} \tag{4.64}
\end{equation*}


ni hosil qilamiz.\\
(4.64) formulaning ikkinchi qoˋshiluvchisini (4.59) ga asosan quyidagicha ifodalab olamiz:


\begin{align*}
& \frac{R}{4 \pi|y|\left|x-y^{*}\right|}=\frac{R}{4 \pi|y|\left|x-y^{*}\right|} \cdot \frac{|y|}{|y|}=\frac{R|y|}{4 \pi|y|^{2}\left|x-y^{*}\right|}= \\
& =\frac{R|y|}{\left.4 \pi|x| y\right|^{2}-y^{*}|y|^{2} \mid}=\frac{R|y|}{\left.4 \pi|x| y\right|^{2}-R^{2} y \mid} \tag{4.65}
\end{align*}


(4.65) ifodani (4.64) formulaga qoˋyib, shar uchun Grin funksiyasini quyidagi ko'rinishda topamiz:


\begin{equation*}
G(x, y)=\frac{1}{4 \pi|x-y|}-\frac{R|y|}{\left.4 \pi|x| y\right|^{2}-R^{2} y \mid} \tag{4,66}
\end{equation*}


2-masala. $\Delta u(x, y, z)=0$ tenglamaning $D_{1}: z>0$ yarim fazodagi Grin funksiyasi $G(x, y, z ; \xi, \eta, \zeta)=\frac{1}{4 \pi}\left(\frac{1}{r}-\frac{1}{r_{1}}\right)$ ko rinishga ega boˋlishini koˋrsating, bu yerda


\begin{align*}
& r=\sqrt{(x-\xi)^{2}+(y-\eta)^{2}+(z-\zeta)^{2}} \\
& r_{1}=\sqrt{(x-\xi)^{2}+(y-\eta)^{2}+(z+\zeta)^{2}} \tag{4.2-chizma}
\end{align*}


Yechish. (4.56) va (4.11) asosan, Grin funksiyasi

$$
G(x, y ; \xi, \eta)=\frac{1}{4 \pi}\left(\frac{1}{r}+g(x, y, z ; \xi, \eta, \zeta)\right)
$$

ko'rinishga ega bo'ladi. $g(x, y, z ; \xi, \eta, \zeta)$ funksiya, qaralayotgan sohada garmonik funksiya boˋlgani uchun $\Delta g(x, y, z)=0$ va $\Delta g(\xi, \eta, \zeta)=0$ tenglamalarni qanoatlantiradi, shu bilan birga (4.57)\\
shartga asosan $(x, y, z)$ ва $(\xi, \eta, \zeta)$ nuqtalarning hech bo'lmaganda bittasi $\partial D_{1}$ ga tegishli bo'lsa, $G(x, y, z, \xi, \eta, \zeta)=0$, ya'ni $g(x, y, z ; \xi, \eta, \zeta)=-1 / r$ bo lishi kerak. Demak, $(x, y, z)$ вa $(\xi, \eta, \zeta)$ nuqtalarning hech bo'lmaganda bittasi $\partial D_{1}$ ga tegishli boˋlishi uchun $z=0$ yoki $\zeta=0$ boˋlishi shart. Bundan esa $r=r_{1}$ boˋladi, ya'ni $g(x, y, z ; \xi, \eta, \zeta)=-1 / r_{1}$. Bu funksiya $D_{1}$ sohada garmonikdir. Shunday qilib,

$$
G(x, y ; \xi, \eta)=\frac{1}{4 \pi}\left(\frac{1}{r}+g(x, y, z ; \xi, \eta, \zeta)\right)=\frac{1}{4 \pi}\left(\frac{1}{r}-\frac{1}{r_{1}}\right)
$$

funksiya $D_{1}$ sohada Laplas tenglamasining Grin funksiyasi bo ladi.\\
$\Delta u(x, y)=0$ tenglamaning $D_{2}: \quad y>0$ yarim tekislikdagi Grin funksiyasi quyidagi ko'rinishda bo'ladi:


\begin{equation*}
G(x, y ; \xi, \eta)=G(\xi, \eta ; x, y)=\frac{1}{2 \pi}\left[\ln \frac{1}{r_{M M_{0}}}-\ln \frac{1}{r_{M M_{1}}}\right] \tag{4.67}
\end{equation*}


bu yerda $\left.\quad \begin{array}{c}r_{M M M_{0}} \\ r_{M M M_{1}}\end{array}\right\}=\sqrt{(x-\xi)^{2}+(y \mp n)^{2}}$\\
(4.3-chizma).

\begin{figure}[h]
\begin{center}
  \includegraphics[alt={},max width=\textwidth]{f9c87acf-24e6-44a5-be24-c530e7441594-05_421_470_1037_66}
\captionsetup{labelformat=empty}
\caption{4.2 - chizma}
\end{center}
\end{figure}

\begin{figure}[h]
\begin{center}
  \includegraphics[alt={},max width=\textwidth]{f9c87acf-24e6-44a5-be24-c530e7441594-05_394_369_1060_616}
\captionsetup{labelformat=empty}
\caption{4.3 - chizma}
\end{center}
\end{figure}

Integral ifoda [11], [15]. Agar $u(x)$ funkasiya $C^{2}(D) \cap C^{1}(\bar{D})$ sinfga tegishli boˋlsa, u holda bu funksiyaning integral ifodasi


\begin{align*}
& u(x)=\frac{1}{2 \pi} \int_{S}\left[\frac{\partial u(\xi)}{\partial n} E(x, \xi)-u(\xi) \frac{\partial E(x, \xi)}{\partial n}\right] d, S-\frac{1}{2 \pi} \int_{D} E(x, \xi) \Delta u(\xi) d \xi= \\
& =\left\{\begin{array}{l}
\frac{1}{2 \pi} \int_{S}\left[\frac{\partial u(\xi)}{\partial n} \ln \frac{1}{r}-u(\xi) \frac{\partial}{\partial n} \ln \frac{1}{r}\right] d_{\xi} S- \\
\quad-\frac{1}{2 \pi} \int_{D} \ln \frac{1}{r} \Delta u(\xi) d \xi, \quad n=2 \\
\frac{1}{(n-2)\left|S_{1}\right|} \int_{S}\left[\frac{1}{r^{n-2}} \frac{\partial u(\xi)}{\partial n}-u(\xi) \frac{\partial}{\partial n}\left(\frac{1}{r^{n-2}}\right)\right] d_{\xi} S- \\
\quad-\int_{D} \frac{\Delta u(\xi)}{r^{n-2}} d \xi, \quad n>2
\end{array}\right. \tag{4.68}
\end{align*}


ko rinishda bo'ladi, bu yerda $S_{1}$ - birlik sfera.\\
Agar $u(x)$ funksiya $D$ sohada garmonik funksiya bo Isa, u holda (4.68) va (4.69) formulalarda $\Delta u=0$ bo'lib, bu garmonik funksiyaning integral ifodasi

\[
u(x)=\left\{\begin{array}{l}
\frac{1}{2 \pi} \int_{S}\left[\frac{\partial u(\xi)}{\partial n} \ln \frac{1}{r}-u(\xi) \frac{\partial}{\partial n} \ln \frac{1}{r}\right] d_{\xi} S, \quad n=2,  \tag{4.70}\\
\frac{1}{(n-2) \mid S_{1} \int_{S}}\left[\frac{1}{r^{n-2}} \frac{\partial u(\xi)}{\partial n}-u(\xi) \frac{\partial}{\partial n}\left(\frac{1}{r^{n-2}}\right)\right] d_{\xi} S, \quad n>2
\end{array}\right.
\]

koˋrinishda boˋladi.\\
Yuqoridagi integral ifodalarga $0^{\circ}$ xshash ushbu $\Delta u(x, y)=-f(x)$ Puasson tenglamasi uchun

$$
\left.u\right|_{S}=\varphi(x), x \in S
$$

Dirixle masalasining yechimi


\begin{equation*}
u(x)=-\int_{S} \frac{\partial G}{\partial n_{y}} \cdot \varphi(y) d S_{y}+\int_{D} G(x, y) f(y) d y \tag{4.72}
\end{equation*}


formula orqali topiladi, bu yerda $G(x, \xi)$ Grin funksiya bo lib, (4.56) formuladan aniqlanadi.

Grin funksiyasi yordamida chegaraviy masalalarni yechish. Yarim tekislikda (4.2) Laplas tenglamasi uchun ( $n=2$ ) Dirixle masalasini qaraymiz:


\begin{gather*}
\Delta u \equiv \Delta u(x, y)=0, \quad y>0, \quad-\infty<x<+\infty,  \tag{4.73}\\
u(x, 0)=f(x), \quad-\infty<x<+\infty . \tag{4.74}
\end{gather*}


(4.73), (4.74) masalaning yechimi quyidagi ko'rinishda


\begin{equation*}
u(x, y)=\left.\int_{-\infty}^{+\infty} f(\xi) \cdot \frac{\partial G(x, y ; \xi, \eta)}{\partial \eta}\right|_{\eta=0} d \xi \tag{4.75}
\end{equation*}


ifodalanadi, bu yerda $G(x, y, \xi, \eta)$ - funksiya (4.67) formula orqali aniqlanadi, ya'ni

$$
G(x, y ; \xi, \eta)=\frac{1}{2 \pi}\left[\ln \frac{1}{\sqrt{(x-\xi)^{2}+(y-\eta)^{2}}}-\ln \frac{1}{\sqrt{(x-\xi)^{2}+(y+\eta)^{2}}}\right]
$$

Bu funksiyani $\eta$ boˋyicha differensiallab, quyidagiga ega boˋlamiz:


\begin{equation*}
\frac{\partial G}{\partial \eta}=-\frac{1}{2 \pi} \cdot \frac{-(y-\eta)}{(x-\xi)^{2}+(y-\eta)^{2}}+\frac{1}{2 \pi} \cdot \frac{(y+\eta)}{(x-\xi)^{2}+(y+\eta)^{2}} . \tag{4.76}
\end{equation*}


(4.76) ni (4.75) ga qo yib, (4.73), (4.74) masalaning yechimini


\begin{equation*}
u(x, y)=\frac{y}{\pi} \int_{-\infty}^{+\infty} \frac{f(\xi)}{(x-\xi)^{2}+y^{2}} d \xi \tag{4.77}
\end{equation*}


olamiz.\\
3-masala. Quyidagi


\begin{gather*}
\Delta u(x, y)=0, \quad y>0, \quad-\infty<x<+\infty,  \tag{4.73}\\
u(x, 0)=\frac{x}{x^{2}+1}, \quad-\infty<x<+\infty . \tag{4.78}
\end{gather*}


Dirixle masalasini yeching.\\
Yechish. (4.78) ni (4.77) ga qo'yib, quyidagini hisoblaymiz:


\begin{equation*}
u(x, y)=\frac{y}{\pi} \int_{-\infty}^{+\infty} \frac{\xi d \xi}{\left[(x-\xi)^{2}+y^{2}\right]\left(1+\xi^{2}\right)}=2 i y[\operatorname{resf}(i)+\operatorname{resf}(x+i y)], \tag{4.79}
\end{equation*}


bu yerda $f(z)=\xi /\left[(x-\xi)^{2}+y^{2}\right]\left(1+\xi^{2}\right)$.\\[0pt]
Bundan, chegirmalar haqidagi teoremadan foydalanib, [9], [14]


\begin{equation*}
\operatorname{res} f(i)=\frac{1}{2\left[(i-x)^{2}+y^{2}\right]}, \quad \operatorname{res} f(x+i y)=\frac{x+i y}{2 i y\left[1-(x+i y)^{2}\right]} \tag{4,80}
\end{equation*}


ni topamiz.\\
Buni (4.79) ga qo yamiz:

$$
u(x, y)=\frac{y}{\pi} \int_{-\infty}^{\infty} \frac{\xi d \xi}{\left[(x-\xi)^{2}+y^{2}\right]\left(1+\xi^{2}\right)}=\frac{i y}{(i-x)^{2}+y^{2}}+\frac{x+i y}{1-(x+i y)^{2}}=
$$

$$
\begin{gathered}
=\frac{i y}{[(i-x)+i y][(i-x)-i y]}+\frac{x+i y}{[(x+i y)-1][(x+i y)+1]}= \\
=\frac{1}{2}\left[\frac{1}{i-x-i y}-\frac{1}{i-x+i y}\right]+\frac{1}{2}\left[\frac{1}{x+i y-1}-\frac{1}{1+x+i y}\right]= \\
=\frac{1}{2}\left[-\frac{1}{i(y-1)+x}-\frac{1}{i(1+y)-x}+\frac{1}{x+i(y-1)}+\frac{1}{x+i(y+1)}\right]= \\
=\frac{1}{2}\left[-\frac{1}{i(1+y)-x}+\frac{1}{x+i(y+1)}\right]=\frac{1}{2}\left[\frac{x-i(y+1)}{(1+y)^{2}+x^{2}}+\frac{x+i(y+1)}{x^{2}+(y+1)^{2}}\right]= \\
=\frac{x}{x^{2}+(1+y)^{2}} .
\end{gathered}
$$

Shunday qilib, (4.73), (4.78) masalaning yechimi

$$
u(x, y)=\frac{x}{x^{2}+(1+y)^{2}}
$$

koˋrinishda boˋladi.\\
Yarim fazoda (4.2) Laplas tenglamasi uchun ( $n=3$ ) Dirixle masalasini qaraymiz:

$$
\begin{gathered}
\Delta u=\Delta u(x, y, z)=0, \quad-\infty<x, y<+\infty, \quad z>0, \\
u(x, y, 0)=f(x, y), \quad-\infty<x, y<+\infty .
\end{gathered}
$$

Bu masalaning yechimi quyidagi ko rinishda

$$
u(x, y, z)=\left.\int_{-\infty}^{+\infty} \int_{-\infty}^{+\infty} f(\xi, \eta) \frac{\partial G}{\partial \zeta}\right|_{\zeta=0} d \xi d \eta
$$

ifodalanadi, bu yerda $G(x, y, z ; \xi, \eta, \zeta)$ - funksiya ushbu formula orqali aniqlanadi ( 4 -§ dagi 2 -masalaga qarang)

$$
G(x, y ; \zeta, \eta)=\frac{1}{4 \pi}\left[\frac{1}{\sqrt{(x-\xi)^{2}+(y-\eta)^{2}+(z-\zeta)^{2}}}-\frac{1}{\sqrt{(x-\xi)^{2}+(y+\eta)^{2}+(z+\zeta)^{2}}}\right] .
$$

Bu funksiyani $\zeta$ boˋyicha differensiallab, so'ng $\zeta=0$ qoyib, quyidagiga ega boˋlamiz:

$$
\left.\frac{\partial G}{\partial \zeta}\right|_{\zeta=0}=\frac{1}{4 \pi} \cdot \frac{2 z}{\left[(x-\xi)^{2}+(y-\eta)^{2}+z^{2}\right]^{3 / 2}} .
$$

Demak, yarim fazoda $(z>0)$ (4.2) Laplas tenglamasi uchun Dirixle masalasining yechimi


\begin{equation*}
u(x, y, z)=\frac{1}{4 \pi} \int_{-\infty}^{+\infty} \int_{-\infty}^{+\infty} \frac{2 z}{\left[(x-\xi)^{2}+(y-\eta)^{2}+z^{2}\right]^{1 / 2}} f(\xi, \eta) d \xi d \eta \tag{4.80}
\end{equation*}


koˋrinishda boˋladi.\\
4-masala. Quyidagi

$$
\begin{gathered}
\Delta u(x, y, z)=0, \quad-\infty<x, y<+\infty, \quad z>0 \\
u(x, y, 0)=\cos x \cdot \cos y, \quad-\infty<x, y<+\infty
\end{gathered}
$$

Dirixle masalasini yeching.\\
Yechish. Qo yilgan Dirixle masalasining yechimi quyidagi

$$
u(x, y, z)=\frac{z}{2 \pi} \int_{-\infty}^{+\infty} \int_{-\infty}^{+\infty} \frac{\cos \xi \cdot \cos \eta d \xi d \eta}{\left[(\xi-x)^{2}+(\eta-y)^{2}+z^{2}\right]^{3 / 2}}
$$

formuladan topiladi. Bu integralni hisoblash uchun $\xi-x=w, \quad \eta-y=v$ almashtirishiami bajarib,


\begin{gather*}
w(x, y, z)=\frac{z}{2 \pi} \int_{-\infty}^{+\infty} \int_{-\infty}^{+\infty} \frac{\cos (x+w) \cdot \cos (v+y) d w d u}{\left[w^{2}+v^{2}+z^{2}\right]^{1 / 2}}= \\
=\frac{z}{2 \pi} \int_{-\infty}^{+\infty} \int_{-\infty}^{+\infty} \frac{\cos w \cdot \cos x-\sin w \cdot \sin x}{\left[w^{2}+v^{2}+z^{2}\right]^{3 / 2}}[\cos v \cdot \cos y-\sin v \cdot \sin y] d w d u= \\
=\frac{z}{2 \pi} \cos x \cdot \cos y \int_{-\infty}^{+\infty} \int_{-\infty}^{+\infty} \frac{\cos w \cdot \cos v d w d u}{\left[w^{2}+v^{2}+z^{2}\right]^{3 / 2}} \tag{4.81}
\end{gather*}


ni hosil qilamiz, ifodadagi qolgan uchta integral, integral ostidagi funksiyalarning toqligiga ko'ra nolga teng boˋladi.

Endi ushbu integralni hisoblaymiz:

$$
\begin{gathered}
J=\int_{-\infty}^{+\infty} \int_{-\infty}^{+\infty} \frac{\cos w \cdot \cos v d w d v}{\left[w^{2}+v^{2}+z^{2}\right]^{3 / 2}}=\iint_{-\infty}^{+\infty} \frac{\cos (w+v)+\sin w \cdot \sin v d w d v}{\left[w^{2}+v^{2}+z^{2}\right]^{3 / 2}}= \\
=\int_{-\infty}^{+\infty} \int_{-\infty}^{+\infty} \frac{\cos (w+v) d v d v}{\left[w^{2}+v^{2}+z^{2}\right]^{3 / 2}}
\end{gathered}
$$

chunki

$$
\int_{-\infty}^{+\infty} \int \frac{\sin w \cdot \sin v d w d v}{\left[w^{2}+v^{2}+z^{2}\right]^{1 / 2}}=0
$$

Oxirgi integralda $\mu=\frac{1}{\sqrt{2}}(w+v), \quad \rho=\frac{1}{\sqrt{2}}(w-v)$ almashtirishlarni bajarib,


\begin{gather*}
J=\int_{-\infty}^{+\infty} \int_{-\infty}^{+\infty} \frac{\cos (\sqrt{2} \mu) d \mu d \rho}{\left[\mu^{2}+\rho^{2}+z^{2}\right]^{\mu 2}}=\int_{-\infty}^{+\infty} \cos (\sqrt{2} \mu) d \mu \int_{-\infty}^{+\infty} \frac{d \rho}{\left[\mu^{2}+\rho^{2}+z^{2}\right]^{\lambda / 2}}= \\
=\left\{\rho=\sqrt{\mu^{2}+z^{2}} \operatorname{tg} t\right\}=\int_{-\infty}^{+\infty} \cos (\sqrt{2} \mu) d \mu \int_{-\pi / 2}^{\pi / 2} \frac{|\cos t| d t}{\mu^{2}+z^{2}}= \\
=2 \int_{-\infty}^{+\infty} \frac{\cos (\sqrt{2} \mu)}{\mu^{2}+z^{2}} d \mu \tag{4.82}
\end{gather*}


ga ega bo'lamiz.\\[0pt]
(4.82) formulaga chegirmalar haqidagi Koshi [14] teoremasini qoˋllaymiz:

$$
\begin{gathered}
J=\operatorname{Re} \int_{-\infty}^{\infty} \frac{\left(e^{i \sqrt{2} \mu}+e^{-i \sqrt{2} \mu}\right)}{\mu^{2}+z^{2}} d \mu=2 \operatorname{Re} \int_{0}^{\infty} \frac{e^{i \sqrt{2} \mu}}{\mu^{2}+z^{2}} d \mu=\left.4 \pi i r e s \frac{e^{i \sqrt{2} \mu}}{\mu^{2}+z^{2}}\right|_{\mu=z i}= \\
=4 \pi i \cdot \frac{e^{-z \sqrt{2}}}{2 z i}=\frac{2 \pi}{z} \cdot e^{-z \sqrt{2}}
\end{gathered}
$$

Bundan va (4.81)dan qo yilgan Dirixle masalasining yechimini

$$
u(x, y, z)=e^{-z \sqrt{2}} \cos x \cdot \cos y
$$

koˋrinishda topamiz.\\
Boshqa usullar yordamida chegaraviy masalalarni yechish\\
1- izoh. (4.73), (4.78) masalaning yechimini Grin funksiyasidan foydalanmay ham topish mumkin. Haqiqatan,

$$
u(x, y)=\frac{1}{2 \pi} \ln \frac{1}{\sqrt{x^{2}+(y+\alpha)^{2}}}(\alpha>0)
$$

funksiya yuqori yarim tekislikda $(y>0)$ (4.73) Laplas tenglamasini qanoatlantiradi, ya'ni

$$
\Delta u(x, y)=\Delta \ln \frac{1}{\sqrt{x^{2}+(y+\alpha)^{2}}}=0
$$

Bundan tashqari, quyidagi tengliklar o'rinli:

$$
\frac{\partial}{\partial x} \Delta \ln \frac{1}{\sqrt{x^{2}+(y+\alpha)^{2}}}=0 \text { yoki } \Delta \frac{\partial}{\partial x} \ln \frac{1}{\sqrt{x^{2}+(y+\alpha)^{2}}}=0
$$

ya'ni

$$
\Delta\left(\frac{x}{r^{2}}\right)=0, \quad r^{2}=x^{2}+(y+\alpha)^{2} .
$$

Shunday qilib, $u(x, y)=\frac{x}{x^{2}+(y+\alpha)^{2}}$ funksiya yuqori yarim tekislikda garmonikdir. Bundan, (4.78) chegaraviy shartga ko'ra $\alpha=1$ ga teng boˋlib, $u(x, y)=\frac{x}{x^{2}+\left(1+y^{\prime}\right)^{2}}$ funksiya (4.73), (4.78) masalaning yechimi bo'ladi.

2-izoh. Agar $z=x+i y, \frac{y}{(x-\xi)^{2}+y^{2}}=\operatorname{Re} \frac{1}{i(\xi-z)}$ boˋlsa, u holda, (4.77) formulani quyidagi ko'rinishda


\begin{equation*}
u(z)=\operatorname{Re} \frac{1}{\pi i} \int_{-\infty}^{+\infty} \frac{u(\xi)}{\xi-z} d \xi \tag{4.83}
\end{equation*}


yozib olamiz.\\
Endi yarim tekislikda (Im $z>0$ ) (ya'ni $y>0$ da) Laplas tenglamasi uchun Dirixle masalasini koˋraylik:

$$
\begin{gathered}
\Delta u(x, y)=0, \quad-\infty<x<+\infty, \quad y>0 \\
\left.u(x, y)\right|_{y=0}=R(x), \quad-\infty<x<+\infty
\end{gathered}
$$

bu yerda $R(z)$ haqiqiy oˋqda qutb maxsus nuqtaga ega boˋlmagan haçiqiy ratsional funksiya bo'lib, $z \rightarrow \infty$ da $R(z) \rightarrow 0$.

Berilgan Dirixle masalasining yechimi


\begin{equation*}
u(z)=\operatorname{Re} \frac{1}{\pi i} \int_{-\infty}^{+\infty} \frac{R(\xi)}{\xi-z} d \xi \tag{4.84}
\end{equation*}


formula orqali topiladi. (4.84) formuladagi integral chegirmalar haqidagi Koshi [9], [14] teoremasi yordamida hisoblanadi:


\begin{equation*}
u(z)=-2 \operatorname{Re} \sum_{\operatorname{Im} \zeta_{0}<0} \operatorname{res} \frac{R(\zeta)}{\zeta \zeta_{0}} \frac{(\zeta)}{\zeta-z} \tag{4.85}
\end{equation*}


5-masala. Quyidagi $\Delta u(x, y)=0, \quad-\infty<x<+\infty, \quad y>0$,

$$
u(x, 0)=\frac{k}{x^{2}+1}, \quad-\infty<x<+\infty, \quad k=\text { const }
$$

Dirixle masalasini yeching.\\
Yechish. (4.84) va (4.85) formulalardan foydalanib, Dirixle masalasining yechimini quyidagi koˋrinishda topamiz:

$$
u(z)=-2 \operatorname{Re} \operatorname{res}_{\zeta=-i} \frac{k}{(\zeta-z)\left(1+\zeta^{2}\right)}=-2 \operatorname{Re} \frac{k}{2 i(z+i)}=\frac{k(y+1)}{x^{2}+(y+1)^{2}}
$$

6-masala. Quyidagi

$$
\begin{gathered}
\Delta u(x, y, z)=z e^{-i} \sin x \cdot \sin y, \quad-\infty<x, y<+\infty, \quad z>0 \\
u(x, y, 0)=0, \quad-\infty<x, y<+\infty
\end{gathered}
$$

masalaning $z \rightarrow+\infty$ da chegaralangan $u(x, y, z)$ yechimini toping.\\
Yechish. Berilgan masala yechimini

$$
u(x, y, z)=\vartheta(z) \cdot \sin x \cdot \sin y
$$

koˋrinishda izlaymiz. U holda

$$
\begin{gathered}
\Delta u \equiv \frac{\partial^{2} u}{\partial x^{2}}+\frac{\partial^{2} u}{\partial y^{2}}+\frac{\partial^{2} u}{\partial z^{2}}=-\vartheta(z) \cdot \sin x \cdot \sin y-\vartheta(z) \cdot \sin x \cdot \sin y+ \\
+\vartheta^{\prime \prime}(z) \cdot \sin x \cdot \sin y
\end{gathered}
$$

ni hisobga olib, quyidagi tenglikka

$$
-2 \vartheta(z) \cdot \sin x \cdot \sin y+\vartheta^{\prime \prime}(z) \cdot \sin x \cdot \sin y=z e^{-z} \sin x \cdot \sin y
$$

ega bo'lamiz. Bundan quyidagi

$$
\vartheta^{\prime \prime}(z)-2 \vartheta(z)=z e^{-z}
$$

oddiy differensial tenglamani hosil qilamiz. Uning umumiy yechimi


\begin{equation*}
\theta(z)=c_{1} e^{\sqrt{2} z}+c_{2} e^{-\sqrt{2} z}+e^{-z}(2-z) \tag{4.86}
\end{equation*}


koˋrinishda boˋladi.\\
Endi $c_{1}$ va $c_{2}$ oˋzgarmaslarni chegaraviy shartlar yordamida topamiz. (4.86) yechimda $c_{1}=0$, chunki agar $c_{1} \neq 0$ bo'lsa, u holda $z \rightarrow+\infty$ da $\vartheta(z) \rightarrow \infty$ bo'ladi. Shuning uchun (4.86) yechimda $z=0 \operatorname{deb}, 0=\vartheta(0)=c_{2}+2 \Rightarrow c_{2}=-2$ aniqlaymiz. Demak, (4.86) ga koˋra qo yilgan masalaning yechimi

$$
u(x, y, z)=\left[e^{-z}(2-z)-2 e^{-\sqrt{2} z}\right] \sin x \cdot \sin y
$$

koˋrinishda topamiz.\\
7-masala. Quyidagi $\Delta u(x, y, z)=0, \quad-\infty<x, y<+\infty, z>0$,

$$
u(x, y, 0)=\frac{x^{2}+y^{2}-2}{\left(x^{2}+y^{2}+1\right)^{5 / 2}}, \quad-\infty<x, y<+\infty
$$

masalaning $u(x, y, z)$ yechimini toping.

Yechish. Shuni aytish mumkinki, $z>0$ yarim fazoning hamma joyida $u(x, y, z)=\frac{1}{\sqrt{x^{2}+y^{2}+(1+z)^{2}}}$ funksiya $u_{x x}+u_{y y}+u_{z z}=0$ Laplas tenglamasini qanoatlantiradi, ya'ni

$$
\Delta\left[\frac{1}{\sqrt{x^{2}+y^{2}+(1+z)^{2}}}\right]=0 .
$$

Bu tenglikni $z$ bo yicha ikki marta differensiallab, quyidagiga


\begin{equation*}
\Delta\left[\frac{x^{2}+y^{2}-2(1+z)^{2}}{\left(x^{2}+y^{2}+(1+z)^{2}\right)^{5 / 2}}\right]=0 \tag{4.87}
\end{equation*}


ega boˋlamiz. Bundan xulosa qilib, shuni aytish lozimki


\begin{equation*}
u(x, y, z)=\frac{x^{2}+y^{2}-2(1+z)^{2}}{\left(x^{2}+y^{2}+(1+z)^{2}\right)^{5 / 2}} \tag{4.88}
\end{equation*}


funksiya $z>0$ yarim fazoning hamma joyida garmonik bo'lib ((4.87) ga qarang), $z=0$ da quyidagi

$$
\left.u\right|_{5=0}=\frac{x^{2}+y^{2}-2}{\left(x^{2}+y^{2}+1\right)^{5 / 2}}
$$

chegaraviy shartni qanoatlantiradi.\\
Shunday qilib, quyilgan masalaning yechimi (4.88) formula orqali topiladi.

Bigarmonik tenglama. Ushbu $\Delta^{2} u=0$ bigarmonik va $\Delta^{2} u=f$ tenglamalar uchun Grin funksiyasidan foydalanmay, chegaraviy masalalarni oˋrganish mumkin.

8-masala. Quyidagi $K=\{(r, \varphi): 0 \leq r \leq a, 0 \leq \varphi<2 \pi\}$ sohada

$$
\begin{gathered}
\Delta^{2} u=0 \quad K \partial a \\
\left.u\right|_{r=a}=0,\left.\quad \frac{\partial u}{\partial n}\right|_{r=a}=A \cos \varphi \quad \partial K \partial a
\end{gathered}
$$

masalani yeching, $n$ - doiraning chegarasida o'tkazilgan tashqi birlik normal.

Yechish. Bizga ma'lumki, bigarmonik tenglama uchun Dirixle masalasi:

$$
\left.u\right|_{r=a}=\psi,\left.\quad \frac{\partial u}{\partial n}\right|_{r=a}=\phi
$$

ushbu ko rinishdagi yechimga ega


\begin{align*}
u(r, \varphi)= & \frac{1}{2 a \pi}\left(r^{2}-a^{2}\right)^{2}\left[-\frac{1}{2} \int_{0}^{2 \pi} \frac{\phi(t) d t}{r^{2}+a^{2}-2 a r \cos (\varphi-t)}+\right. \\
& \left.+\int_{0}^{2 \pi} \frac{\psi(t)[a-r \cos (\varphi-t)] d t}{\left[r^{2}+a^{2}-2 a r \cos (\varphi-t)\right]^{2}}\right] \tag{4.89}
\end{align*}


(4.89) formulaga $\psi=0, \quad \phi=A \cos \varphi$ qo'yib,

$$
w(r, \varphi)=-\frac{A}{2 a}\left(r^{2}-a^{2}\right)^{2} \cdot \int_{0}^{2 \pi} \frac{1}{2 \pi} \cdot \frac{\cos t d t}{r^{2}+a^{2}-2 a r \cos (\varphi-t)}
$$

ni hosil qilamiz. IV bobning 3 -§iga koˋra oxirgi integral Laplas tenglamasi uchun qo'yilgan

$$
\begin{gathered}
\Delta v=0 \quad K \partial a \\
v_{t=a}=\cos \varphi \quad \partial K \partial a
\end{gathered}
$$

Dirixle masalasining yechimidir. Bu masalaning yechimi $v=r / a \cos \varphi$ koˋrinishda boˋladi. Demak, qoˋyilgan masala yechimining yagonaligiga ko'ra qo'yidagi tenglikka

$$
\frac{1}{2 \pi} \cdot \int_{0}^{2 \pi} \frac{\cos t d t}{r^{2}+a^{2}-2 a r \cos (\varphi-t)}=\frac{r}{a} \cos \varphi
$$

ega boˋlamiz.\\
Shunday qilib, bigarmonik tenglama uchun qo'yilgan Dirixle masalasining yechimi

$$
u(r, \varphi)=-\frac{\operatorname{Ar}\left(r^{2}-a^{2}\right)^{2}}{2 a^{2}} \cos \varphi
$$

koˋrinishda boˋladi.\\
9-masala. Quyidagi $K=\{(r, \varphi): 0 \leq r \leq a, 0 \leq \varphi<2 \pi\}$ sohada

$$
\begin{gathered}
\Delta^{2} u=1 \quad K \partial a, \\
\left.u\right|_{r=a}=0,\left.\quad \frac{\partial u}{\partial n}\right|_{r=a}=0, \quad \partial K \partial a
\end{gathered}
$$

masalani yeching, $n$ - doiraning chegarasada o'tkazilgan tashqi birlik\\
normal.\\
Yechish. Faraz qilaylik, izlanayotgan yechim faqat $r$ ga bogˋliq bo'lsin, ya'ni $u=u(r)$. U holda

$$
\Delta^{2} u=\left(\frac{\partial^{2}}{\partial r^{2}}+\frac{\partial}{r \partial r}\right)^{2} u=\left(\frac{\partial^{4}}{\partial r^{4}}+\frac{2 \partial^{3}}{r \partial r^{3}}+\frac{\partial^{2}}{r^{2} \partial r^{2}}\right) u
$$

oˋrinlidir.\\
Shunday qilib, oddiy differensial tenglama uchun quyidagi masalaga kelamiz:

$$
\frac{d^{4} u}{d r^{4}}+\frac{2 d^{3} u}{r d r^{3}}+\frac{d^{2} u}{r^{2} d r^{2}}=1,0<r<a,\left.\quad u\right|_{r=a}=0,\left.\quad \frac{\partial u}{\partial r}\right|_{r=a}=0
$$

Bu masalani yechib, qo'yilgan masalaning yechimini

$$
u(r)=\frac{a^{4}}{7}\left[\frac{1}{12}\left(\frac{r}{a}\right)^{4}-\frac{1}{3}\left(\frac{r}{a}\right)+\frac{1}{4}\right]
$$

ko rinishda topamiz.\\
10-masala. Quyidagi $E=\{(x, y):-\infty<x<+\infty, y>0\}$ yarim tekislikda $\Delta^{2} u=\left.e^{-2 y} \sin x \quad E \partial a \quad u\right|_{y=0}=0,\left.\frac{\partial u}{\partial y}\right|_{y=0}=0,-\infty<x<+\infty$ masalani yeching.

Yechish. Berilgan tenglamani quyidagi ko rinishda yozib olamiz:

$$
\frac{\partial^{4} u}{\partial x^{4}}+2 \frac{\partial^{4} u}{\partial x^{2} \partial y^{2}}+\frac{\partial^{4} u}{\partial y^{4}}=e^{-2 y} \sin x .
$$

Bu tenglamaning yechimini $u(x, y)=\omega(y) \sin x$ ko'rinishda izlaymiz, $\quad \omega(y)$-noma'lum funksiya. $\quad u(x, y)=\omega(y) \sin x$ funksiyani berilgan tenglamaga qoˋyib,

$$
\omega^{(V)}-2 \omega^{\prime \prime}+\omega=e^{-2 y}
$$

oddiy differensial tenglamani olamiz. Uning umumiy yechimi

$$
\omega(y)=C_{1} e^{y}+C_{2} y e^{y}+C_{3} e^{-y}+C_{4} y e^{-y}+\frac{1}{9} e^{-2 y}
$$

koˋrinishda boˋladi. Masala shartiga koˋra $C_{1}=0, C_{2}=0$ teng, aksincha, $\quad y \rightarrow+\infty \quad$ da $\quad \omega(y) \rightarrow \infty \quad$ bo'ladi. $\quad C_{3}, \quad C_{4} \quad$ lar $\omega(0)=\omega^{\prime}(0)=0$ shartlardan topiladi. Demak, oddiy differensial\\
tenglamaning yechimi

$$
\omega(y)=\frac{1}{9}\left(-e^{-y}+y e^{-y}+e^{-2 y}\right)
$$

ko'rinishda bo'lib, qo'yilgan masalasining yechimi esa

$$
u(x, y)=\frac{1}{9}\left(-e^{-y}+y e^{-y}+e^{-2 y}\right) \sin x
$$

formuladan topiladi.

\section*{Mustaqil yechish uchun masalalar}
\begin{enumerate}
  \setcounter{enumi}{592}
  \item $\Delta u(x, y, z)=0$ tenglamaning $y>0, z>0$ ikki yoqli\\
burchakdagi Grin funksiyasi
\end{enumerate}

$$
G(x, y, z ; \xi, \eta, \zeta)=\frac{1}{4 \pi}\left(\frac{1}{r}-\frac{1}{r_{1}}-\frac{1}{r_{2}}+\frac{1}{r_{3}}\right)
$$

koˋrinishga ega bolishini koˋrsating. Bu yerda

$$
\begin{array}{ll}
r=\sqrt{(x-\xi)^{2}+(y-\eta)^{2}+(z-\zeta)^{2}}, & r_{1}=\sqrt{(x-\zeta)^{2}+(y-\eta)^{2}+(z+\zeta)^{2}} \\
r_{2}=\sqrt{(x-\xi)^{2}+(y+\eta)^{2}+(z-\zeta)^{2}}, & r_{3}=\sqrt{(x-\zeta)^{2}+(y+\eta)^{2}+(z+\zeta)^{2}}
\end{array}
$$

\begin{enumerate}
  \setcounter{enumi}{593}
  \item $\Delta u(x, y, z)=0$ tenglamaning $x>0, y>0, z>0$ oktantadagi
\end{enumerate}

Grin funksiyasi

$$
G(x, y, z ; \xi, \eta, \zeta)=\frac{1}{4 \pi} \sum_{k=0}^{7}(-1)^{k} \frac{1}{r_{k}}
$$

koˋrinishga ega bo'lishini ko'rsating. Bu yerda

$$
\begin{array}{ll}
r_{0}=\sqrt{(x-\xi)^{2}+(y-\eta)^{2}+(z-\zeta)^{2}}, & r_{1}=\sqrt{(x-\xi)^{2}+(y-\eta)^{2}+(z+\zeta)^{2}}, \\
r_{2}=\sqrt{(x-\xi)^{2}+(y+\eta)^{2}+(z+\zeta)^{2}}, & r_{3}=\sqrt{(x-\xi)^{2}+(y+\eta)^{2}+(z-\zeta)^{2}}, \\
r_{4}=\sqrt{(x+\xi)^{2}+(y+\eta)^{2}+(z-\zeta)^{2}}, & r_{5}=\sqrt{(x+\xi)^{2}+(y-\eta)^{2}+(z-\zeta)^{2}}, \\
r_{6}=\sqrt{(x+\zeta)^{2}+(y-\eta)^{2}+(z+\zeta)^{2}}, & r_{7}=\sqrt{(x+\xi)^{2}+(y+\eta)^{2}+(z+\zeta)^{2}} .
\end{array}
$$

\begin{enumerate}
  \setcounter{enumi}{594}
  \item $\Delta u=0$ tenglamaning radiusi $R=4$ ga teng boˋlgan shardagi
\end{enumerate}

Grin funksiyasi

$$
G(x, y)=\frac{1}{4 \pi}\left(\frac{1}{|x-y|}-\frac{4}{|y|\left|x-y^{*}\right|}\right)
$$

koˋrinishga ega boˋlishini koˋrsating.\\
596. $\Delta u(x, y, z)=0$ tenglamaning ikkita $z=0$ va $z=l$ tekisliklar bilan chegaralangan sohadagi Grin funksiyasi

$$
G(x, y, z ; \xi, \eta, \zeta)=\frac{1}{4 \pi} \sum_{n=-\infty}^{\infty}\left(\frac{1}{r_{n}}-\frac{1}{r_{n}^{\prime}}\right)
$$

koˋrinishga ega boˋlishini koˋrsating, bu yerda

$$
\begin{aligned}
& r_{n}=\sqrt{(x-\xi)^{2}+(y-\eta)^{2}+[z-(2 n l+\zeta)]^{2}} \\
& r_{n}^{\prime}=\sqrt{(x-\xi)^{2}+(y-\eta)^{2}+[z-(2 n l-\zeta)]^{2}}
\end{aligned}
$$

\begin{enumerate}
  \setcounter{enumi}{596}
  \item $\Delta u(x, y, z)=0$ tenglamaning $z>0$ yarim fazoda $z=0$ da quyidagi $\frac{\partial u}{\partial z}+h u=0$ chegaraviy shartni qanoatlantiradigan Grin funksiyasini toping.
  \item $\Delta u(x, y, z)=0$ tenglamaning quyidagi sohadagi Grin funksiyasini toping:
\end{enumerate}

\begin{enumerate}
  \item $x^{2}+y^{2}+z^{2}<R^{2}, \quad z>0$ yarim sharda;
  \item $x^{2}+y^{2}+z^{2}<R^{2}, \quad y>0 \quad z>0$ sharning toˋrtdan bir qismida:
  \item $x^{2}+y^{2}+z^{2}<R^{2}, x>0 \quad y>0 \quad z>0$ sharning sakkizdan bir qismida.
\end{enumerate}

\begin{enumerate}
  \setcounter{enumi}{598}
  \item $\Delta u(z)=0, z=x+i y$ tenglamaning quyidagi sohadagi Grin funksiyasini toping:
\end{enumerate}

\begin{enumerate}
  \item Im $z>0$ yarim tekislikda; 2) $0<\arg z<0,5 \pi$ tekislikning to rtdan birida;
  \item $|z|<R$ doirada;'
  \item $|z|<R \quad \operatorname{Im} z>0$ yarim doirada;
  \item $|z|<1,0<\arg z<0,5 \pi$ doiraning to'rtdan birida;
  \item $0<\operatorname{Im} z<\pi$ "polasa"da;
  \item $0<\operatorname{Im} z<\pi, \operatorname{Re} z>0$ yarim polasada;
\end{enumerate}

\begin{enumerate}
  \setcounter{enumi}{599}
  \item $D=\left\{(x, y):-1<x<1,0<y<\sqrt{1-x^{2}}\right\}$ sohada
\end{enumerate}

$$
\left\{\begin{array}{l}
\Delta u(x, y)=0, \quad-1<x<1, \quad 0<y<\sqrt{1-x^{2}} \\
\left.u\right|_{x^{2}+y^{2}=r^{2}}=\varphi(s), \quad 0 \leq s \leq l,\left.\quad \frac{\partial u}{\partial y}\right|_{y=0}=v(x)-1<x<1
\end{array}\right.
$$

masalaning Grin funksiyasini tuzing va yechimni toping.\\
601. $D=\left\{(x, y):-1<x<1,0<y<\sqrt{1-x^{2}}\right\}$ sohada\\
$\left\{\begin{array}{l}\Delta u(x, y)=0, \quad-1<x<1, \quad 0<y<\sqrt{1-x^{2}}, \\ \left.u\right|_{x^{2}+y^{2}=r^{2}}=\varphi(s), \quad 0 \leq s \leq l,\left.\quad u\right|_{y=0}=\tau(x), \quad-1 \leq x \leq 1\end{array}\right.$\\
masalaning Grin funksiyasini tuzing va yechimni toping.\\
602. $D=\left\{(x, y):-1<x<1,0<y<\left[\left(\frac{m+2}{2}\right)^{2}\left(1-x^{2}\right)\right]^{1 /(m+2)}\right\}$ sohada

$$
\left\{\begin{array}{l}
y^{m} u_{x x}+u_{y y}=0,\left.\quad u\right|_{\Gamma}=\varphi(s), \quad 0 \leq s \leq l \\
\left.u\right|_{y=0}=\tau(x), \quad-1 \leq x \leq 1 ; \quad \Gamma: x^{2}+\frac{4}{(m+2)^{2}} y^{m+2}=1
\end{array}\right.
$$

masalaning Grin funksiyasini tuzing va yechimni toping.\\
603. $D=\left\{(x, y):-1<x<1,0<y<\left[\left(\frac{m+2}{2}\right)^{2}\left(1-x^{2}\right)\right]^{1 /(m+2)}\right\}$ sohada

$$
\left\{\begin{array}{l}
y^{m} u_{x x}+u_{y y}=0,\left.\quad u\right|_{\Gamma}=\varphi(s), \quad 0 \leq s \leq l, \\
\left.\frac{\partial u}{\partial y}\right|_{y=0}=v(x), \quad-1<x<1 ; \quad \Gamma: x^{2}+\frac{4}{(m+2)^{2}} y^{m+2}=1
\end{array}\right.
$$

masalaning Grin funksiyasini tuzing va yechimini toping.\\
604. Ushbu $\Delta u=f(x, y, z),-\infty<x, y<+\infty, z>0,\left.u\right|_{z=0}=u_{0}(x, y)$,\\
$-\infty<x, y<+\infty$ Dirixle masalasini $f$ va $u_{0}$ quyidagi qiymatlarida yeching. ( $f$ va $u_{0}$ - bo'lakli-uzluksiz va chegaralangan funksiyalar)

\begin{enumerate}
  \item $f=e^{-2} \sin x \cdot \cos y, u_{0}=0$;
  \item $f=0, \quad u_{0}=\theta(y-x)= \begin{cases}1, & y-x \geq 0, \\ 0, & y-x<0 ;\end{cases}$
  \item $f=0, \quad u_{0}=\left(1+x^{2}+y^{2}\right)^{-1 / 2}$;
  \item $f=2\left((1+z)^{2}+x^{2}+y^{2}\right)^{-2}, \quad u_{0}=\left(1+x^{2}+y^{2}\right)^{-1}$;
\end{enumerate}

\begin{enumerate}
  \setcounter{enumi}{604}
  \item Ushbu $\Delta u(x)=0, \quad x_{2}>0, \quad x_{3}>0, \quad u_{s}=u_{0}(x), \quad x=\left(x_{1}, x_{2}, x_{3}\right)$
\end{enumerate}

Dirixle masalasini $u_{0}$ ning quyidagi qiymatlarida yeching:

\begin{enumerate}
  \item $u_{0}$ - boˋlakli-uzluksiz va chegaralangan funksiya;
  \item $\left.u_{0}\right|_{x_{2}=0}=0,\left.\quad u_{0}\right|_{x_{3}=0}=\exp \left\{-4 x_{1}\right\} \cdot \sin 5 x_{2}$;
  \item $\left.u_{0}\right|_{x_{2}=0}=0,\left.\quad u_{0}\right|_{x_{3}=0}=x_{2}\left(1+x_{1}^{2}+x_{2}^{2}\right)^{-1 / 2}$.
\end{enumerate}

\begin{enumerate}
  \setcounter{enumi}{605}
  \item Dirixle masalasini $|x|<R$ sharda yeching:
\end{enumerate}

$$
\Delta u(x)=-f(x), \quad|x|<R,\left.u\right|_{|x|=R}=\varphi(x), \quad|x|=\sqrt{x_{1}^{2}+x_{2}^{2}+x_{3}^{2}}
$$


\end{document}